\chapter{The Curvature of Spherically Symmetric Spacetimes}
\label{ch:sphmetrics}

\begin{usesbox}
\textbf{This chapter uses:} the curvature of warped products (\cref{prop:warped}); the frame formalism (\cref{ch:frames}).\\
\textbf{It provides:} the Christoffel symbols, Ricci tensor, scalar curvature and Einstein tensor of a general spherically symmetric metric in two gauges, the diagonal gauge $-N^2\dd t^2+A^2\dd x^2+R^2\dd\Omega^2$ and the double-null gauge $-2\Omega\,\dd u\,\dd v+r^2\dd\Omega^2$, all derived from the Levi-Civita formula; and the Gauss--Codazzi data of a spherically symmetric spatial slice.
\end{usesbox}

Every spherically symmetric spacetime can be written as a warped product of a two-dimensional Lorentzian metric $h_{ab}$ and the round sphere,
\[
\dd s^2=h_{ab}\,\dd x^a\dd x^b+r^2\,\dd\Omega^2,\qquad \dd\Omega^2=\dd\theta^2+\sin^2\theta\,\dd\phi^2,
\]
where $r$, the areal radius, is a function on the two-dimensional quotient. The Einstein tensor of such a metric is the input to everything in Part~IV. Most texts quote it; here it is derived, in two gauges that will both be needed, from the coordinate formulas
\[
\Gamma^a{}_{bc}=\tfrac12g^{ae}(\partial_bg_{ec}+\partial_cg_{eb}-\partial_eg_{bc}),\qquad
R_{bc}=\partial_a\Gamma^a{}_{bc}-\partial_c\Gamma^a{}_{ba}+\Gamma^a{}_{ae}\Gamma^e{}_{bc}-\Gamma^a{}_{ce}\Gamma^e{}_{ba},
\]
$R=g^{bc}R_{bc}$ and $G_{bc}=R_{bc}-\frac12g_{bc}R$. As in \cref{ch:frames} the derivation is algebraic: functions form a commutative algebra with four commuting derivations, $\partial_\theta\sin\theta=\cos\theta$ and $\partial_\theta\cos\theta=-\sin\theta$. We write subscripts for partial derivatives, $R_t=\partial_tR$ and so on.

\section{The diagonal gauge}

Consider
\begin{equation}\label{eq:diagmetric}
\dd s^2=-N^2\dd t^2+A^2\dd x^2+R^2\dd\Omega^2,
\end{equation}
with $N,A,R$ nonvanishing functions of $(t,x)$. Here $N$ is the lapse, $A\,\dd x$ the proper radial distance and $R$ the areal radius.

\begin{theorem}[Einstein tensor in the diagonal gauge]\label{thm:einD}\ver\
For \eqref{eq:diagmetric}, the inverse metric is diagonal, the Christoffel symbols are given by an explicit table, and
\begin{align*}
\frac{G_{tt}}{N^2}&=\frac{1}{R^2}\Bigl(1-\frac{R_x^2}{A^2}\Bigr)-\frac{2}{A^2R}\Bigl(R_{xx}-\frac{A_xR_x}{A}\Bigr)+\frac1{N^2}\Bigl(\frac{R_t^2}{R^2}+\frac{2A_tR_t}{AR}\Bigr),\\
G_{tx}&=\frac2R\Bigl(\frac{A_t}{A}R_x+\frac{N_x}{N}R_t-R_{tx}\Bigr),\\
\frac{G_{xx}}{A^2}&=-\frac{1}{R^2}\Bigl(1-\frac{R_x^2}{A^2}\Bigr)+\frac{2N_xR_x}{A^2NR}-\frac1{N^2}\Bigl(\frac{R_t^2}{R^2}+\frac{2R_{tt}}{R}-\frac{2N_tR_t}{NR}\Bigr),
\end{align*}
together with explicit expressions for $R_{tt},R_{xx},R_{\theta\theta},R_{\phi\phi},R_{tx}$, the scalar curvature and $G_{\theta\theta}=G_{\phi\phi}/\sin^2\theta$.
\end{theorem}
\begin{proof}
\emph{Christoffel symbols.} For a diagonal metric the formula reduces to $\Gamma^a{}_{aa}=\partial_ag_{aa}/(2g_{aa})$, $\Gamma^a{}_{ab}=\partial_bg_{aa}/(2g_{aa})$ and $\Gamma^a{}_{bb}=-\partial_ag_{bb}/(2g_{aa})$ for $a\ne b$. With $g=\operatorname{diag}(-N^2,A^2,R^2,R^2\sin^2\theta)$ the non-zero symbols are
\begin{align*}
&\Gamma^t{}_{tt}=\tfrac{N_t}N,\quad \Gamma^t{}_{tx}=\tfrac{N_x}N,\quad \Gamma^t{}_{xx}=\tfrac{AA_t}{N^2},\\
&\Gamma^x{}_{tt}=\tfrac{NN_x}{A^2},\quad \Gamma^x{}_{tx}=\tfrac{A_t}A,\quad \Gamma^x{}_{xx}=\tfrac{A_x}A,\\
&\Gamma^t{}_{\theta\theta}=\tfrac{RR_t}{N^2},\quad \Gamma^x{}_{\theta\theta}=-\tfrac{RR_x}{A^2},\quad \Gamma^\theta{}_{t\theta}=\Gamma^\phi{}_{t\phi}=\tfrac{R_t}R,\\
&\Gamma^\theta{}_{x\theta}=\Gamma^\phi{}_{x\phi}=\tfrac{R_x}R,\quad \Gamma^t{}_{\phi\phi}=\sin^2\theta\,\Gamma^t{}_{\theta\theta},\quad \Gamma^x{}_{\phi\phi}=\sin^2\theta\,\Gamma^x{}_{\theta\theta},\\
&\Gamma^\theta{}_{\phi\phi}=-\sin\theta\cos\theta,\quad \Gamma^\phi{}_{\theta\phi}=\cot\theta .
\end{align*}

\emph{A covariant shortcut.} Inserting these symbols into the Ricci formula, and organising the result with the two-dimensional covariant derivative $\nabla$ of $h=-N^2\dd t^2+A^2\dd x^2$, gives the standard formulae for any metric $h_{ab}\dd x^a\dd x^b+R^2\dd\Omega^2$ \cite{MTW}:
\[
G_{ab}=-\frac2R\bigl(\nabla_a\nabla_bR-h_{ab}\Box R\bigr)-\frac{h_{ab}}{R^2}\bigl(1-(\nabla R)^2\bigr),\qquad
G^\theta{}_\theta=G^\phi{}_\phi=\frac{\Box R}R-\frac12\,{}^{(2)}\!R .
\]
(One may check them on Schwarzschild, $h=-f\dd t^2+\dd r^2/f$ with $f=1-2M/r$: then $\Box r=f'$, $(\nabla r)^2=f$, ${}^{(2)}\!R=-f''$, and all components vanish.)

\emph{The diagonal gauge.} With the symbols above,
\[
\begin{aligned}
&\nabla_t\nabla_tR=R_{tt}-\tfrac{N_t}NR_t-\tfrac{NN_x}{A^2}R_x,\qquad \nabla_x\nabla_xR=R_{xx}-\tfrac{AA_t}{N^2}R_t-\tfrac{A_x}AR_x,\\
&\nabla_t\nabla_xR=R_{tx}-\tfrac{N_x}NR_t-\tfrac{A_t}AR_x,
\end{aligned}
\]
$\Box R=-N^{-2}\nabla_t\nabla_tR+A^{-2}\nabla_x\nabla_xR$ and $(\nabla R)^2=-R_t^2/N^2+R_x^2/A^2$. Then
\[
\nabla_t\nabla_tR+N^2\Box R=\frac{N^2}{A^2}\Bigl(R_{xx}-\frac{A_xR_x}A\Bigr)-\frac{A_tR_t}A,
\]
and $G_{tt}=-\frac2R(\nabla_t\nabla_tR+N^2\Box R)+\frac{N^2}{R^2}(1-(\nabla R)^2)$ gives the stated $G_{tt}/N^2$. Similarly $\nabla_x\nabla_xR-A^2\Box R=\frac{A^2}{N^2}(R_{tt}-\frac{N_t}NR_t)-\frac{N_x}NR_x$ gives $G_{xx}/A^2$, and $G_{tx}=-\frac2R\nabla_t\nabla_xR$ gives the mixed component.

For the Schwarzschild check: $N=A^{-1}=\sqrt f$, $R=x$, no $t$-dependence. Then $G_{tx}=0$, and $G_{tt}/N^2=\frac{1-f}{x^2}+\frac{2A_x}{A^3x}=\frac{1-f}{x^2}-\frac{f'}x=\frac{2M}{x^3}-\frac{2M}{x^3}=0$; $G_{xx}$ vanishes in the same way.
\end{proof}


Two special cases are worth reading off at once. On a time-symmetric slice, where all $t$-derivatives vanish, $G_{tt}/N^2$ equals one half of the scalar curvature of the spatial metric $A^2\dd x^2+R^2\dd\Omega^2$. In proper-distance gauge $A=1$ this is
\begin{equation}\label{eq:3R}
\tfrac12\,{}^{(3)}R=\frac{1-R'^2}{R^2}-\frac{2R''}{R},
\end{equation}
which is \cref{prop:warped} again: the three sectional curvatures $-R''/R$, $-R''/R$ and $(1-R'^2)/R^2$ add up to $\frac12\,{}^{(3)}R$. And for $N=A^{-1}=\sqrt{1-2Gm/R}$, $R=x$, with constant $m$, one finds $G_{tt}=G_{xx}=G_{tx}=0$: the Schwarzschild metric.

\section{The double-null gauge}

For dynamical problems it is often better to use null coordinates,
\begin{equation}\label{eq:nullmetric}
\dd s^2=-2\Omega\,\dd u\,\dd v+r^2\dd\Omega^2 .
\end{equation}

\begin{theorem}[Einstein tensor in the double-null gauge]\label{thm:einN}\ver\
For \eqref{eq:nullmetric},
\begin{align*}
G_{uu}&=-\frac2r\Bigl(r_{uu}-\frac{\Omega_u}{\Omega}r_u\Bigr),\qquad
G_{vv}=-\frac2r\Bigl(r_{vv}-\frac{\Omega_v}{\Omega}r_v\Bigr),\\
G_{uv}&=\frac{\Omega}{r^2}+\frac{2r_ur_v}{r^2}+\frac{2r_{uv}}{r},\\
G_{\theta\theta}&=r^2\Bigl(\frac{\Omega_u\Omega_v}{\Omega^3}-\frac{\Omega_{uv}}{\Omega^2}\Bigr)-\frac{2r\,r_{uv}}{\Omega},
\end{align*}
with the corresponding Christoffel symbols, Ricci components (including $R_{uv}=R_{vu}$) and scalar curvature.
\end{theorem}
\begin{proof}
For $h=-2\Omega\,\dd u\,\dd v$ the only non-zero two-dimensional Christoffel symbols are $\Gamma^u{}_{uu}=\Omega_u/\Omega$ and $\Gamma^v{}_{vv}=\Omega_v/\Omega$, and $h^{uv}=-1/\Omega$. Hence
\[
\nabla_u\nabla_ur=r_{uu}-\frac{\Omega_u}\Omega r_u,\qquad \nabla_u\nabla_vr=r_{uv},\qquad \Box r=2h^{uv}r_{uv}=-\frac{2r_{uv}}\Omega,\qquad (\nabla r)^2=-\frac{2r_ur_v}\Omega .
\]
The formulae of the previous proof give $G_{uu}=-\frac2r\nabla_u\nabla_ur$, the same for $v$, and
\[
G_{uv}=-\frac2r\bigl(r_{uv}-(-\Omega)\Box r\bigr)+\frac\Omega{r^2}\bigl(1-(\nabla r)^2\bigr)=\frac{2r_{uv}}r+\frac\Omega{r^2}+\frac{2r_ur_v}{r^2}.
\]
For the angular component we need the curvature of $h$. Writing $\Omega=e^{2\sigma}$, $h$ is conformal to the flat metric $-2\,\dd u\,\dd v$, and ${}^{(2)}\!R=4e^{-2\sigma}\sigma_{uv}=\frac2\Omega\bigl(\frac{\Omega_{uv}}\Omega-\frac{\Omega_u\Omega_v}{\Omega^2}\bigr)$. Then $G_{\theta\theta}=r^2\bigl(\Box r/r-\frac12{}^{(2)}\!R\bigr)$ is the stated expression.
\end{proof}


The equations $G_{uu}=8\pi GT_{uu}$ and $G_{vv}=8\pi GT_{vv}$ are the Raychaudhuri equations for the two families of radial null geodesics; with $\Omega=e^{2\sigma}$ they read $r_{uu}-2\sigma_ur_u=-4\pi Gr\,T_{uu}$. They play the central role in the analysis of null junctions in Chapter~17 and of collapse in Chapter~20.

\section{Spherical slices: Gauss and Codazzi}

Let a spacelike slice carry the metric $\dd s^2+r(s)^2\dd\Omega^2$ and extrinsic curvature $K^i{}_j=\operatorname{diag}(a,b,b)$, the most general spherically symmetric form. The Einstein equations split into the Hamiltonian constraint ${}^{(3)}R+K^2-K_{ij}K^{ij}=16\pi G\rho$ and the momentum constraint $\nabla_jK^j{}_i-\partial_iK=8\pi G\,j_i$.

\begin{proposition}[Codazzi and Gauss data]\label{prop:codazzi}\ver\
With Christoffel symbols computed from $\dd s^2+r^2\dd\Omega^2$:
\begin{enumerate}[label=(\roman*)]
\item $\nabla_jK^j{}_s-\partial_sK=2\frac{r'}{r}(a-b)-2b'$, and the angular components of $\nabla_jK^j{}_i-\partial_iK$ vanish.
\item $K^2-K^i{}_jK^j{}_i=4ab+2b^2$.
\end{enumerate}
\end{proposition}
\begin{proof}
(i) For $\dd s^2+r^2\dd\Omega^2$ the relevant Christoffel symbols are $\Gamma^\theta{}_{s\theta}=\Gamma^\phi{}_{s\phi}=r'/r$, $\Gamma^s{}_{\theta\theta}=-rr'$, $\Gamma^s{}_{\phi\phi}=-rr'\sin^2\theta$ and $\Gamma^\phi{}_{\theta\phi}=\cot\theta$. For the $s$-component,
\[
\nabla_jK^j{}_s=\partial_sK^s{}_s+\Gamma^j{}_{js}K^s{}_s-\Gamma^k{}_{js}K^j{}_k=a'+\frac{2r'}ra-\frac{2r'}rb,
\]
and $\partial_sK=a'+2b'$; subtracting gives $2\frac{r'}r(a-b)-2b'$. For the $\theta$-component, $\partial_\theta K^\theta{}_\theta=0$, and the two Christoffel terms are both $\cot\theta\,b$ and cancel; $\partial_\theta K=0$. The $\phi$-component is similar.

(ii) $K=a+2b$ and $K^i{}_jK^j{}_i=a^2+2b^2$, so $K^2-K^i{}_jK^j{}_i=(a+2b)^2-a^2-2b^2=4ab+2b^2$.
\end{proof}


So in spherical symmetry the momentum constraint is the single equation $2(r'/r)(a-b)-2b'=8\pi G\,j_s$, and the Hamiltonian constraint is ${}^{(3)}R+4ab+2b^2=16\pi G\rho$ with ${}^{(3)}R$ from \eqref{eq:3R}.

\begin{litbox}
These formulas are standard and appear, in various notations, in every treatment of spherical collapse; see for instance Misner and Sharp \cite{MisnerSharp1964}, Poisson \cite{PoissonToolkit} and the double-null formalism used by Christodoulou \cite{Christodoulou1993}. We include them, with complete verified derivations from the Levi-Civita formula, because everything in Part~IV depends on them and a sign error here would propagate everywhere. No new claim is made.
\end{litbox}

\section*{What this chapter passes on}
\begin{itemize}[nosep]
\item The Einstein tensor in the diagonal gauge, and ${}^{(3)}R$ of a spherical slice (\cref{thm:einD}, \eqref{eq:3R}).
\item The Einstein tensor in the double-null gauge (\cref{thm:einN}).
\item The spherical momentum and Hamiltonian constraints (\cref{prop:codazzi}).
\end{itemize}

\section*{Exercises}
\begin{exercise}
Using \cref{thm:einD}, show that $N=A^{-1}=\sqrt{1-2Gm/x}$, $R=x$ gives $G_{tt}=0$.
\end{exercise}
\begin{exercise}
Using \eqref{eq:3R}, compute ${}^{(3)}R$ for $R(s)=a\sin(s/a)$ and for $R(s)=a$ (a cylinder $\RR\times S^2$).
\end{exercise}
\begin{exercise}
Check \cref{prop:codazzi}(ii) directly from $K=a+2b$ and $K^i{}_jK^j{}_i=a^2+2b^2$.
\end{exercise}

\begin{checklist}
\chk{AlphaS4_Curv:giD_inv}{The inverse of the diagonal metric.$^*$}
\chk{AlphaS4_Curv:gamD}{Christoffel symbols of the diagonal metric (table).$^*$}
\chk{AlphaS4_Curv:ricD_00,AlphaS4_Curv:ricD_11,AlphaS4_Curv:ricD_22,AlphaS4_Curv:ricD_33,AlphaS4_Curv:ricD_01}{Ricci components $R_{tt},R_{xx},R_{\theta\theta},R_{\phi\phi},R_{tx}$ (five statements).$^*$}
\chk{AlphaS4_Curv:rscD}{Scalar curvature of the diagonal metric.$^*$}
\chk{AlphaS4_Curv:einD_00}{$G_{tt}$.$^*$}
\chk{AlphaS4_Curv:einD_11}{$G_{xx}$.$^*$}
\chk{AlphaS4_Curv:einD_22}{$G_{\theta\theta}$.$^*$}
\chk{AlphaS4_Curv:einD_01}{$G_{tx}$.$^*$}
\chk{AlphaS4_Curv:giN_inv}{The inverse of the double-null metric.$^*$}
\chk{AlphaS4_Curv:gamN}{Christoffel symbols of the double-null metric.$^*$}
\chk{AlphaS4_Curv:ricN_00,AlphaS4_Curv:ricN_11,AlphaS4_Curv:ricN_22,AlphaS4_Curv:ricN_33,AlphaS4_Curv:ricN_01,AlphaS4_Curv:ricN_10}{Ricci components in the double-null gauge, including $R_{uv}=R_{vu}$ (six statements).$^*$}
\chk{AlphaS4_Curv:rscN}{Scalar curvature in the double-null gauge.$^*$}
\chk{AlphaS4_Curv:einN_00}{$G_{uu}=-\frac2r(r_{uu}-\frac{\Omega_u}{\Omega}r_u)$.$^*$}
\chk{AlphaS4_Curv:einN_11}{$G_{vv}$.$^*$}
\chk{AlphaS4_Curv:einN_01}{$G_{uv}$.$^*$}
\chk{AlphaS4_Curv:einN_22}{$G_{\theta\theta}$ in the double-null gauge.$^*$}
\chk{AlphaIII_Codazzi:codazzi_s}{Radial component of the momentum constraint, $2\frac{r'}{r}(a-b)-2b'$.$^*$}
\chk{AlphaIII_Codazzi:codazzi_theta,AlphaIII_Codazzi:codazzi_phi}{Angular components vanish (two statements).$^*$}
\chk{AlphaIII_Codazzi:hamiltonian_K}{$K^2-K^i{}_jK^j{}_i=4ab+2b^2$.$^*$}
\end{checklist}
